\subsection{Sprint 3}

A Sprint 3 concentrou um conjunto mais amplo de funcionalidades e melhorias na plataforma. As atividades previstas incluíam o aprimoramento da tela de detalhe dos projetos, o dashboard de indicadores de impacto da ONG, a edição de projetos, a user story de visibilidade e acesso ao perfil da ONG, além do upload e visualização de documentos institucionais. Todas as atividades previstas foram concluídas, o que representou um avanço importante na maturidade do produto.

Essa sprint exigiu maior articulação entre frontend, backend, validações de negócio e experiência do usuário. Diferentemente de tarefas isoladas de interface, as funcionalidades previstas dependiam de dados reais, regras de acesso, integração com autenticação e comportamento coerente entre diferentes páginas. Por isso, minha atuação foi se deslocando cada vez mais para atividades de análise, apoio técnico, revisão e organização do trabalho do squad.

No início da sprint, dediquei tempo ao estudo do estado atual do backend e à criação de subtasks referentes a essa camada. Esse trabalho foi importante para transformar demandas amplas em tarefas menores e mais executáveis. Também participei de weeklys e sugeri a criação de uma nova tabela no banco de dados, a partir das necessidades percebidas durante o desenvolvimento. Além disso, revisei pull requests relacionados à edição de ONGs, alinhei decisões com integrantes da AGES III do squad e busquei entender melhor a construção dos processos de tratamento e processamento de documentos no backend.

Uma das atividades mais relevantes foi auxiliar um colega do squad na configuração da geração de tokens do Auth0 no Postman. Esse tipo de apoio foi importante porque a autenticação era uma parte central do projeto e, sem conseguir simular corretamente os perfis e permissões, o desenvolvimento e os testes ficavam travados. Esse momento reforçou a importância de dominar não apenas o código da funcionalidade, mas também o ambiente necessário para que os colegas consigam trabalhar.

Mais adiante, atuei na revisão de pull requests, no apoio a colegas com erros e na implementação da user story de listagem e download de documentos. Essa funcionalidade dialogava diretamente com uma necessidade importante das ONGs: armazenar e disponibilizar documentos institucionais de forma organizada e segura. Ao trabalhar nessa frente, precisei considerar aspectos de frontend, endpoints, permissões e experiência de uso, principalmente para garantir que os documentos fossem exibidos de forma compreensível e pudessem ser acessados quando necessário.

Apesar das entregas terem sido concluídas, o principal problema encontrado foi a gestão de pessoas. A sprint deixou claro que, conforme o projeto cresce, o desafio deixa de ser apenas técnico. A coordenação entre pessoas, expectativas, ritmos de trabalho e responsabilidades passa a ocupar um papel central. Em alguns momentos, percebi dificuldade em manter todos alinhados, tanto dentro do squad quanto com pessoas de fora dele.

A principal lição aprendida foi a necessidade de melhorar o canal de comunicação entre os membros do squad e também com outros squads. A falta de comunicação clara não gera apenas atraso; ela cria insegurança, sobreposição de tarefas e perda de contexto. Para um projeto com múltiplos perfis, muitas funcionalidades e diferentes níveis de experiência entre os integrantes, esse alinhamento precisa ser constante.

Ao final da sprint, o squad ficou com tarefas de refatoração e limpeza do código. Esse próximo passo foi importante porque, após várias entregas funcionais, tornou-se necessário melhorar a organização interna da aplicação, reduzir inconsistências e preparar o projeto para as etapas finais. Para mim, a Sprint 3 foi uma fase de aprendizado mais voltada à sustentação de um projeto em crescimento, especialmente na relação entre apoio técnico, revisão de código, comunicação e gestão de pessoas.
